\chapter{FN-DSA (Falcon): Draft-Oriented Walkthrough}
\chaptermark{FN-DSA (Falcon)}
\label{ch:fndsa}

\section{Learning objectives}

This chapter is a draft-oriented walkthrough of FN-DSA, NIST's planned standardization of Falcon~\cite{fips206,falcon2020}. It is not a claim of conformance to a final FIPS: readers must consult the exact draft version and date used by an implementation. Unlike ML-KEM, ML-DSA, and SLH-DSA, FIPS 206 remains under standardization at this book's standards snapshot. Where ML-DSA uses Fiat--Shamir with aborts, Falcon takes a different route on the same lattice foundation: \emph{hash-and-sign} over NTRU lattices, using a Gaussian trapdoor sampler. The payoff is remarkably compact signatures and keys; the price is the most delicate implementation of the three.

After finishing this chapter, you should be able to:

\begin{enumerate}
  \item contrast hash-and-sign with the Fiat--Shamir approach of ML-DSA;
  \item describe the NTRU lattice and why a short secret basis is a trapdoor;
  \item explain how signing is a nearest-lattice-point problem solved with a Gaussian sampler;
  \item connect the sampler back to discrete Gaussians, good vs.\ bad bases, and CVP from Part~\ref{part:lattices};
  \item understand why Falcon's floating-point arithmetic makes constant-time implementation hard.
\end{enumerate}

\section{Review: hash-and-sign versus Fiat--Shamir}

ML-DSA built a signature from an interactive identification protocol collapsed by a hash. Falcon uses the older, more direct \emph{hash-and-sign} paradigm, the same shape as classical RSA signatures:

\begin{enumerate}
  \item hash the message to a target point $c$;
  \item use the secret trapdoor to produce a short vector that ``explains'' $c$;
  \item the verifier checks the vector is short and consistent with $c$.
\end{enumerate}

The difficulty that stalled lattice hash-and-sign for years was that a naive short vector \emph{leaks the secret basis}: every signature is a hint about the trapdoor. Falcon's core contribution is a sampler whose output distribution is provably independent of the secret -- a discrete Gaussian over the lattice, in the hash-and-sign framework of Gentry, Peikert, and Vaikuntanathan~\cite{gpv2008} --- \textbf{GPV}, after its authors, and the name this chapter uses for it from here on. This is the same ``the mask must hide the secret'' principle as ML-DSA's aborts, achieved by a different mechanism.

\section{The NTRU lattice and its trapdoor}

Falcon rests on the NTRU assumption introduced in Section~\ref{sec:ntru}. It works in the by-now-familiar ring $R_q = \mathbb{Z}_q[x]/(x^n+1)$ with $n \in \{512, 1024\}$ and $q = 12289$. Key generation picks two \emph{short} secret polynomials $f, g \in R_q$ and publishes
\[
h = g \cdot f^{-1} \pmod q.
\]
The public key is just the single polynomial $h$. It defines the \textbf{NTRU lattice}
\[
\Lambda = \{ (u, v) \in R_q^2 : u + v\,h = 0 \pmod q \},
\]
a $2n$-dimensional lattice exactly of the kind studied in Part~\ref{part:lattices}. From $h$ alone you can write down a basis for $\Lambda$, but it is a \emph{bad} basis -- long and skewed -- so finding short vectors in it is hard.

The secret is a \emph{good} basis for the very same lattice. Key generation completes $(f,g)$ into a full short basis by solving the \textbf{NTRU equation}
\[
f\,G - g\,F = q
\]
for two more short polynomials $F, G$. The pair of rows $(g, -f)$ and $(G, -F)$ is a short, nearly orthogonal basis of $\Lambda$ -- precisely the ``good basis'' that Chapters~\ref{ch:lattice}--\ref{ch:svp} told us is the key to solving lattice problems.

\begin{figure}[H]
\centering
\begin{tikzpicture}[scale=0.5]
  \draw[pqcgray!40, thin] (-6.5,0) -- (6.5,0);
  \draw[pqcgray!40, thin] (0,-5.5) -- (0,5.5);
  \foreach \i in {-2,...,2} {
    \foreach \j in {-2,...,2} {
      \pgfmathsetmacro{\x}{2*\i+1*\j}
      \pgfmathsetmacro{\y}{1*\i+3*\j}
      \node[latticept] at (\x,\y) {};
    }
  }
  \node[latticeptbold] at (0,0) {};
  % target point c (off-lattice)
  \node[circle, fill=pqcorange, inner sep=0pt, minimum size=5pt,
        label={[font=\scriptsize,pqcorange]above right:{$c$ (hashed message)}}] (c) at (3.2,1.6) {};
  % nearest lattice point
  \node[latticeptbold, fill=pqcteal] (near) at (3,1) {};
  \draw[shortvec] (near) -- (c);
  \node[pqcteal, font=\scriptsize, below=1pt of near] {closest lattice point};
\end{tikzpicture}
\caption{Signing as a closest-vector problem. Hashing the message gives a target point $c$ (orange). The signer uses the secret \emph{good} basis of the NTRU lattice to find a lattice point very close to $c$; the short difference vector is the signature. With only the public \emph{bad} basis (from $h$) this closest point cannot be found -- that gap is the trapdoor.}
\label{fig:falcon-cvp}
\end{figure}

\section{Phase 1: key generation}

\begin{algorithm}[H]
\caption{FN-DSA.KeyGen (teaching form)}
\label{alg:fndsa-keygen}
\begin{algorithmic}[1]
\Require Ring $R_q$, $q=12289$, dimension $n$
\Ensure Public key $h$; secret key (short basis) $(f,g,F,G)$
\Repeat
  \State sample short $f, g \in R_q$ (small Gaussian coefficients)
\Until{$f$ is invertible mod $q$ \textbf{and} $(f,g)$ is short enough}
\State solve the NTRU equation $fG - gF = q$ for short $F, G$
\State $h \gets g \cdot f^{-1} \bmod q$
\State \Return public $h$, secret $(f,g,F,G)$
\end{algorithmic}
\end{algorithm}

The secret basis and the public $h$ describe the \emph{same} lattice; they differ only in quality. That is the entire trapdoor.

\section{Phase 2: hash the message to a point}

Signing begins by mapping the message to a random-looking target in the ring. A fresh salt $r$ is prepended so that signing the same message twice gives independent targets:
\[
c = \textsc{HashToPoint}(r \,\|\, M) \in R_q.
\]

\section{Phase 3: sample a short solution (the trapdoor sampler)}
\label{sec:falcon-tree-note}

The signer must find a short pair $(s_1, s_2)$ satisfying
\[
s_1 + s_2\,h = c \pmod q,
\]
which is the same as finding a lattice point of $\Lambda$ close to $(c, 0)$. Using the good basis, the signer runs \textbf{fast Fourier sampling}: a discrete Gaussian sampler over the lattice coset, implemented efficiently over the FFT representation of the basis (organized as an ``LDL tree''). Two properties matter:

\begin{itemize}
  \item the result is \emph{short}, because the good basis makes nearby lattice points reachable;
  \item its distribution is a \emph{discrete Gaussian} centered at the target, which -- crucially -- is independent of \emph{which} good basis produced it. This is what stops each signature from leaking the secret.
\end{itemize}

This is where the whole book converges: the discrete Gaussian of Chapter~\ref{ch:prob}, the good-versus-bad-basis geometry of Chapters~\ref{ch:lattice}--\ref{ch:lll}, and the closest-vector view of the LWE chapter all appear at once, in service of one signature.

\begin{figure}[H]
\centering
\begin{tikzpicture}[node distance=6mm and 10mm, every node/.style={font=\scriptsize}]
  \node[flownode] (m) {message $M$};
  \node[flownode, right=of m] (hash) {$c=\textsc{HashToPoint}(r\|M)$};
  \node[flownode, right=of hash] (samp) {Gaussian sampler\\[-2pt](uses secret basis)};
  \node[flownode, below=of samp] (s) {short $(s_1,s_2)$\\[-2pt]$s_1+s_2 h = c$};
  \node[flownode, left=of s] (comp) {compress $s_2$};
  \node[keynode, left=of comp] (sig) {signature\\[-2pt]$\sigma=(r, s_2)$};
  \draw[flowarrow] (m) -- (hash);
  \draw[flowarrow] (hash) -- (samp);
  \draw[flowarrow] (samp) -- (s);
  \draw[flowarrow] (s) -- (comp);
  \draw[flowarrow] (comp) -- (sig);
\end{tikzpicture}
\caption{The Falcon signing flow. The message is hashed to a target $c$; the secret basis drives a Gaussian sampler that returns a short solution of $s_1+s_2h=c$; only the salt $r$ and the compressed second component $s_2$ are output, since the verifier can recover $s_1$.}
\label{fig:falcon-sign}
\end{figure}

\section{Phase 4: compress and output}

Only $s_2$ needs to be sent: the verifier can recover $s_1 = c - s_2 h \bmod q$. Because $s_2$ has small Gaussian coefficients, it compresses well, giving Falcon its signature-size advantage. The signature is $\sigma = (r, \text{compress}(s_2))$.

\begin{algorithm}[H]
\caption{FN-DSA.Sign / Verify (teaching form)}
\label{alg:fndsa-signverify}
\begin{algorithmic}[1]
\Statex \textbf{Sign}$(sk=(f,g,F,G),\ M)$:
\State sample salt $r$;\quad $c \gets \textsc{HashToPoint}(r\|M)$
\State $(s_1,s_2) \gets \textsc{FFSampling}$ of a short solution to $s_1+s_2h=c$ using the secret basis
\State \Return $\sigma=(r,\ \textsc{Compress}(s_2))$
\Statex
\Statex \textbf{Verify}$(pk=h,\ M,\ \sigma=(r,\hat s_2))$:
\State $s_2 \gets \textsc{Decompress}(\hat s_2)$;\quad $c \gets \textsc{HashToPoint}(r\|M)$
\State $s_1 \gets c - s_2\,h \bmod q$ \Comment{recover the first component}
\State \Return \emph{accept} if $\|(s_1,s_2)\|^2 \le \beta^2$, else \emph{reject} \Comment{must be short}
\end{algorithmic}
\end{algorithm}

Verification is cheap: one multiplication to recover $s_1$, then a norm check. A forger without the good basis would have to solve the closest-vector problem in the NTRU lattice using only $h$ -- exactly the hardness Part~\ref{part:lattices} established.

\section{Why Falcon is the hardest to implement}
\label{sec:falcon-hard}

The Gaussian sampler needs high-precision \emph{floating-point} arithmetic (FFT over the reals) to get the output distribution exactly right. Floating-point operations are notoriously variable-time and platform-dependent, which makes a \emph{constant-time} implementation -- one that leaks nothing through timing -- genuinely difficult. If the sampler's timing or rounding depends on the secret basis, an attacker can gradually reconstruct it. This is the direct motivation for the final chapter on side-channel security, and it is the main reason Falcon, despite its elegance and compactness, is the most demanding of the three standards to deploy safely.

\subsection{Non-determinism is a second, separate problem}

Timing is the familiar hazard, and it is not the only one. The same floating-point
operation can give slightly different results on a different processor, compiler,
or optimization level, and that alone can break the scheme --- with no side
channel involved at all.

The reason is the security requirement underneath the GPV framework: a signer must
never release \emph{two distinct short preimages for the same digest}. Two short
vectors $s, s'$ with $As = As' = c$ give their difference $s - s'$, a short
non-zero vector in the lattice $\Lambda^{\perp}_q(A)$ --- which is a solution to
the very problem the trapdoor was supposed to keep hard. Ordinary Falcon avoids
this by drawing a fresh 320-bit salt for every signature, so no two signatures
ever target the same digest.

The danger appears when signing is \emph{derandomized}: if the salt is derived
pseudorandomly from the message and the key, as one would want for reproducible
signing, then signing the same message twice must produce the same signature
bit-for-bit. Two implementations at different precision --- or the same
implementation compiled twice --- would produce two valid signatures for one
digest, and hand an observer a short lattice vector.

The fix is to remove the platform dependence rather than to manage it: replace
hardware floating-point with an \emph{integer emulation}, so that every
implementation performs identical arithmetic and produces identical output. The
cost is asymmetric in a way that happens to suit a blockchain
(Section~\ref{sec:bitcoin-scheme-choice}):

\begin{itemize}
  \item key generation is about $2\times$ slower;
  \item signing is about $15\times$ slower;
  \item \textbf{verification is unaffected}, because Falcon verification uses no
        floating-point arithmetic at all --- it is integer arithmetic modulo $q$
        with an NTT, and it is the fastest verification of the three standards.
\end{itemize}

An output is signed once, by its owner, when it is spent; the signature is then
verified by every validating node, forever. Paying $15\times$ on the operation
that happens once to make the operation that happens everywhere both cheap and
identical across platforms is the right side of that trade. Even emulated,
Falcon signing remains one to two orders of magnitude faster than SLH-DSA at its
small-signature parameters.

\subsection{The signer needs memory, too}

One further practical constraint, easy to miss because it appears in neither the
key nor the signature size. The sampler is driven by the tree of
Section~\ref{sec:falcon-tree-note}, and where that tree is kept decides the
signer's memory footprint. Holding it alongside the compact basis makes signing a
simple traversal but occupies roughly $40$ kB at $n=512$ and $90$ kB at
$n=1024$. Rebuilding only the branch currently needed cuts this to about $8$ and
$16$ kB, at roughly double the signing cost.

That matters because the natural place to run a signer is a hardware wallet, and
those are microcontrollers with tight RAM budgets. Note also that this cost
compounds with the emulation penalty above, and that it is a \emph{per-signature}
cost whenever a key signs only once --- which is exactly the recommended practice
for a Bitcoin address.

\section{SageMath experiment: the public NTRU relation}

Sage can verify the public-key equation underlying the NTRU lattice. The
example deliberately stops before the Gaussian trapdoor sampler: finding a
correctly distributed short signature is the hard, specification-sensitive part
of FN-DSA and is not reproduced by this toy calculation.

\begin{lstlisting}[language=Python]
from random import Random
from sage.all import GF, PolynomialRing

q = 17
n = 8
rng = Random("fndsa-ntru-relation")
R.<x> = PolynomialRing(GF(q))
Rq = R.quotient(x^n + 1)

def small_poly():
    return Rq([rng.choice([-1, 0, 1]) for _ in range(n)])

def small_unit():
    f = small_poly()
    while not f.is_unit():
        f = small_poly()
    return f

f = small_unit()
g = small_poly()
h = g * f^(-1)                 # public NTRU value
assert f * h == g               # f h = g in Rq

c = Rq([rng.randrange(q) for _ in range(n)])
s2 = small_poly()
s1 = c - s2*h                  # verifier's reconstruction rule
assert s1 + s2*h == c
print("NTRU public relation and verification equation hold")
\end{lstlisting}

The last equality is only an algebra check. In a real signature, \textsc{FFSampling}
must choose $(s_1,s_2)$ with the required short Gaussian distribution; choosing
an arbitrary small $s_2$ as above is not an FN-DSA signing algorithm.

\section{A complete SageMath implementation}
\label{sec:fips206-impl}

The experiment above stops exactly where the difficulty starts. The
companion file \sagefile{fips206\_fndsa.sage} goes the rest of the way:
all eighteen numbered algorithms of the Falcon specification, at $n=512$ and
$n=1024$ -- the FFT over the complex tower, \textsc{NTRUGen} with the
tower-recursive NTRU solver, the LDL\textsuperscript{$\star$} decomposition
and the Falcon tree, fast Fourier sampling with the isochronous discrete
Gaussian sampler, signing, verification, and the compressed signature and key
encodings.

\playground{the sampler and FFT checks, and signing and verification from a stored key; only \textsc{NTRUGen} is too slow for a sandbox}

\paragraph{Which specification the code follows.}
This matters more here than in the other three chapters. At the book's
standards snapshot NIST lists FIPS~206 as \emph{in development}: no initial
public draft has been published, and consequently no ACVP test vectors exist.
The authoritative specification for FN-DSA is therefore still the round-3
Falcon submission, and that is what the file implements, using Falcon's
algorithm numbering. The places most likely to be respecified when the draft
appears -- the key and signature encodings, and the domain separation of the
hashed message -- are marked in the source. The lattice machinery below them
is not expected to change.

\paragraph{Solving the NTRU equation.}
Key generation's one genuinely hard step is the line that
Algorithm~\ref{alg:fndsa-keygen} states in five words: \emph{solve
$fG - gF = q$ for short $F,G$}. The method is a tower recursion. The field
norm $N$ maps $\mathbb{Z}[x]/(x^n+1)$ down to
$\mathbb{Z}[x]/(x^{n/2}+1)$; applying it repeatedly reaches $\mathbb{Z}$,
where the equation is an extended gcd; the multiplicativity of the norm then
lifts the solution back up, one level at a time.

\begin{lstlisting}[language=Python]
def NTRUSolve(f, g):
    """Algorithm 6. Find F, G with f*G - g*F = q in Z[x]/(x^n + 1)."""
    n = len(f)
    if n == 1:
        d, u, v = xgcd(ZZ(int(f[0])), ZZ(int(g[0])))
        if d != 1:
            return None, None
        return [int(-FALCON_q*v)], [int(FALCON_q*u)]
    fp, gp = field_norm(f), field_norm(g)
    Fp, Gp = NTRUSolve(fp, gp)              # recurse in the smaller ring
    if Fp is None:
        return None, None
    F = _cyclo_mul(lift(Fp), galois_conjugate(g), n)
    G = _cyclo_mul(lift(Gp), galois_conjugate(f), n)
    return Reduce(f, g, F, G)               # shorten F, G against f, g
\end{lstlisting}

Here \texttt{field\_norm} is $N(a) = a(x)a(-x)$ written in $x^2$,
\texttt{lift} is $a(x)\mapsto a(x^2)$, and \texttt{galois\_conjugate} is
$a(x)\mapsto a(-x)$. The lifted $F,G$ come back with enormous coefficients,
which is what \textsc{Reduce} is for: it repeatedly subtracts the multiple of
$(f,g)$ nearest to $(F,G)$, rescaling to 53 significant bits each round so
that the rounding stays exact enough to converge.

\paragraph{Three kinds of arithmetic, three Sage rings.}
Falcon is unusual in needing exact and inexact arithmetic side by side, and
the implementation keeps them in separate, explicitly named Sage domains:

\begin{itemize}
  \item the NTRU solver runs over \texttt{ZZ['x']}, so the trapdoor equation
        $fG-gF=q$ is established \emph{exactly} rather than up to floating
        point -- \texttt{verify\_ntru\_equation()}, at the end of
        \sagefile{fips206\_fndsa.sage}, checks it as an identity
        in $\mathbb{Z}[x]/(x^n+1)$;
  \item the public key and the verification relation
        $s_1 = c - s_2 h$ are computed in
        $\mathbb{F}_{12289}[x]/(x^n+1)$;
  \item the signing FFT runs over \texttt{CDF}, Sage's complex double field,
        which \emph{is} the IEEE~754 binary64 arithmetic that Falcon
        prescribes -- not a higher-precision stand-in that would silently
        change the sampler's output distribution.
\end{itemize}

The function \texttt{verify\_fft\_against\_exact\_ring()}, in the same file,
closes the loop
between the first and third of these, checking that pointwise multiplication
in the FFT domain reproduces the product Sage computes exactly in
$\mathbb{Z}[x]/(x^n+1)$.

\paragraph{Validation without NIST vectors.}
Since FIPS~206 has none, FN-DSA is validated against the Falcon
submission's own known-answer tests~\cite{falcon2020}, in four independent
ways:

\begin{enumerate}
  \item \textsc{SamplerZ} reproduces all sixteen test vectors of the Falcon
        specification's sampler table, byte for byte -- including the cases
        where rejection sampling iterates twice or \textsc{BerExp} consumes a
        second random byte;
  \item for each Falcon KAT key, the private key is decoded, $fG-gF=q$ is
        verified exactly, and the recomputed $h=g f^{-1}$ re-encodes to the
        KAT public key byte for byte;
  \item every Falcon KAT signature verifies, and the same signature over a
        tampered message is rejected;
  \item signatures produced by this implementation under a Falcon KAT
        \emph{private} key verify under the corresponding KAT public key.
\end{enumerate}

Signing itself cannot be pinned to the KAT byte string, and the reason is
instructive: Falcon signing is randomized, and the reference KATs were
generated with the reference implementation's own PRNG chain, so identical
signatures would test PRNG replication rather than the scheme. Verification
is deterministic, and it is checked byte-exactly. The gap is precisely the
Gaussian sampler discussed in the next section -- the part that is hardest to
specify, hardest to implement, and hardest to test.

\section{Parameter sets}

\begin{table}[H]
\centering
\begin{tabular}{lcccc}
\toprule
Scheme & $n$ & NIST level & Public key & Signature \\
\midrule
FN-DSA-512  & 512  & 1 & $\approx 897$ B & $\approx 666$ B \\
FN-DSA-1024 & 1024 & 5 & $\approx 1793$ B & $\approx 1280$ B \\
\bottomrule
\end{tabular}
\caption{The two FN-DSA (Falcon) parameter sets, both with $q=12289$. Signatures and keys are markedly smaller than ML-DSA's or SLH-DSA's -- the reason Falcon exists.}
\end{table}

\paragraph{Why there is no Level 3.}
The gap in that table is structural, not an oversight. The fast Fourier sampler
recurses by halving the ring degree at every step down to a scalar leaf, so $n$
must be a power of two: the admissible degrees are $\dots, 256, 512, 1024,
\dots$ with nothing in between. Concrete estimates place $n=512$ at roughly $128$
bits of classical security and $n=1024$ well above $250$, and no power-of-two
ring lands in the middle. ML-DSA, whose security is tuned by the module
dimensions $(k,\ell)$ over a fixed ring, has no such constraint --- which is why
it offers three levels and Falcon offers two. A deployment wanting a Category 3
margin from a lattice signature must take Falcon-1024 and the size that comes
with it, or take ML-DSA instead.

\section{Pitfalls}

\paragraph{Pitfall 1: sampling short vectors deterministically.}
If the signer always returned the \emph{single} closest lattice point (e.g.\ via Babai rounding), every signature would trace the shape of the secret basis, and a few signatures would reveal it. The Gaussian \emph{randomization} of the output is essential, not a mere convenience.

\paragraph{Pitfall 2: treating the floating-point sampler as ordinary code.}
The sampler's timing, branches, and rounding must not depend on secret data. Naive floating-point code leaks; Falcon implementations go to great lengths (emulated fixed-precision, careful table lookups) to be constant-time.

\paragraph{Pitfall 3: reusing the salt.}
The per-signature salt $r$ must be fresh. Reusing it across different messages correlates their targets $c$ and erodes the Gaussian's hiding property.

\section{Summary}

FN-DSA (Falcon) is lattice hash-and-sign done right:

\begin{itemize}
  \item the public key $h$ and the secret short basis $(f,g,F,G)$ describe the same NTRU lattice, differing only in quality -- the trapdoor;
  \item signing hashes the message to a target and uses a Gaussian sampler over the good basis to find a short, secret-hiding solution;
  \item verification recovers $s_1=c-s_2 h$ and checks shortness, a closest-vector problem for anyone without the trapdoor;
  \item the reward is very compact signatures; the cost is delicate floating-point, constant-time engineering.
\end{itemize}

This completes the \emph{lattice-based} signatures. ML-DSA and FN-DSA both live on lattices but reach a signature by opposite routes -- Fiat--Shamir with aborts versus Gaussian hash-and-sign -- and trade off size against implementation delicacy accordingly. The next part leaves lattices entirely for signatures built on a different hard problem: hash functions alone, culminating in the third signature standard, SLH-DSA (FIPS 205).
